\chapter{What index-calculus hardness rests on}\label{ch:hardness}

Chapter~\ref{ch:falldegree} studied one quantity, the solving degree of the
point-decomposition system, and showed that its first fall degree does not
bound it. A complexity claim for elliptic-curve index calculus needs more
than a degree. It needs a bounded solver cost, a decomposition yield with a
known distribution, relations of sufficient rank, a usable geometry, an
interface that no generic algorithm reproduces, and an end-to-end
composition of all of these. Each of the subexponential claims in the
literature discharges some of these obligations by proof and the rest by an
assumption. This chapter records, assumption by assumption, what the
programme proved, refuted, measured or left open between 2 and
4~November~2026.

The work falls into four kinds:
\begin{itemize}
\item a typed obligation chain and a reanalysis of committed evidence
  (\art{research/ic_weak_assumptions_20261102/});
\item exact measurements of six weak points on Weil-descent systems
  (\art{research/ic_weakpoints_20261102/});
\item exact censuses against five named heuristics
  (\art{research/ecdlp_weakness_census_20261104/}) and Lean~4 theorems
  aimed at the same nodes (\art{research/ecdlp_weakness_lean_20261104/},
  \art{research/lean_attacks_20261102/},
  \art{research/inside_gap_n4_20261102/});
\item a claim graph of the hardness literature and a triage of its weak
  points (\art{docs/ecdlp-hardness-corpus/WEAKNESS_TRIAGE.md}).
\end{itemize}

\paragraph{Standing.} None of this work is an ecbench session, a leaderboard
row or a validated result. No validator report or replay certificate exists
for any census or measurement below. Every measured number in this chapter
is therefore an unverified assertion of its producer under the
proof-carrying-claims rule; the census README marks its numbers
\texttt{[pcc:unverified]}. The Lean results are stated with their exact
formal status in \S\ref{sec:hd-lean}: ``kernel-checked'' is used only where
a Lean kernel accepted the declaration and its source contains no
\texttt{sorry}, \texttt{admit}, \texttt{native\_decide} or user axiom. The
corpus records every internal theorem here as
\texttt{proved\_internal\_unreviewed}: no external reviewer has read the
statements against the literature they address. Nothing in this chapter
changes a row of Tables~A--C, and none of it is a measured improvement of the
index-calculus baseline.

\section{Six serial obligations}\label{sec:hd-chain}

The weak-assumption programme of 2~November
(\art{research/ic_weak_assumptions_20261102/README.md}; literature pass and
reanalysis of committed records, no new run) writes an index-calculus
security argument as a chain of six obligations. The chain is serial in the
sense that matters: failure of any one edge invalidates the complexity
claim, whatever the others show. Its machine-readable form,
\texttt{obligations.json}, is checked by \texttt{audit.py}, a fail-closed
gate that refuses to promote a route while any field is missing. It covers
naive prime-field Semaev, generic pair tables, Kummer (PKM) towers, the
proposed arity-three hybrid solver, Wagner filters, the hashed-base tensor
route (\S\ref{sec:pt-tensor}) and xedni. The route-by-route reading for prime
fields is in \S\ref{sec:pg-obligations}; Table~\ref{tab:hd-chain} lists the
obligations with the proxy usually offered for each and the section of this
chapter that bears on it.

\begin{table}[tbp]
\centering\small
\begin{tabularx}{\linewidth}{@{}>{\raggedright\arraybackslash}p{0.16\linewidth}>{\raggedright\arraybackslash}X>{\raggedright\arraybackslash}p{0.22\linewidth}>{\raggedright\arraybackslash}p{0.2\linewidth}@{}}
\toprule
Obligation & What it needs & Proxy often used instead & Here \\
\midrule
Solution geometry & dimension and radicality of the full system, rational lifts, exceptional charts, target strata & \texttt{zero\_dim=true} on a random target & \S\ref{sec:hd-weakpoints} (W5) \\
\addlinespace
Solver growth & maximal matrix dimensions, nonzeros, arithmetic and memory, with timeouts and out-of-memory exits kept & first fall degree; a stable solving degree & \S\ref{sec:hd-chain}, \S\ref{sec:hd-pq2}, Chapter~\ref{ch:falldegree} \\
\addlinespace
Yield and variance & rational decompositions per natural target: distribution, zero mass, multiplicity, solver failures & the mean count $B^m/r$ & \S\ref{sec:hd-yield} \\
\addlinespace
Relation rank & rank growth, dead columns, full-rank probability & raw row count & \S\ref{sec:hd-rank} \\
\addlinespace
Non-generic interface & a coordinate predicate whose useful candidates no generic table reproduces at equal cost & ``uses coordinates'' & \S\ref{sec:hd-generic} \\
\addlinespace
End-to-end composition & one parameter family charging set-up, failed queries, collection, verification, linear algebra, descent, memory & product of the separate bounds & \S\ref{sec:hd-weakpoints} (W6) \\
\bottomrule
\end{tabularx}
\caption[]{The six obligations of an index-calculus complexity claim, after
\art{research/ic_weak_assumptions_20261102/README.md}. The README's table
also states, per obligation, the ``strongest next falsifier''; the right
column points to the evidence in this chapter.}
\label{tab:hd-chain}
\end{table}

\paragraph{A stable degree does not certify cost.} The reanalysis extracted
the committed PKM $m=4$ random-target F4 rows, two per cell
(\art{research/ic_weak_assumptions_20261102/evidence-summary.json}, key
\texttt{pkm\_m4\_random\_f4}; \texttt{extract\_existing.py --check}
recomputes it from the JSONL records):
\begin{itemize}
\item at $N=8,12,16$ the solving degrees are $6,7,7$;
\item the maximal column counts are $1{,}049$, $9{,}452$ and $76{,}837$;
\item the maximal multiply-add counts are $48{,}447{,}128$,
  $85{,}477{,}602{,}933$ and $27{,}048{,}206{,}854{,}510$, that is
  $4.84\cdot10^7$, $8.55\cdot10^{10}$ and $2.70\cdot10^{13}$.
\end{itemize}
From $N=8$ to $N=16$ the columns grow $73.25$-fold and the arithmetic
$558{,}304$-fold, although the last two degrees are equal. Three cells are
not an asymptotic fit. They are a direct counterexample to reading a
bounded degree as a bounded cost: the cost of a Gröbner computation at a
fixed degree $d$ is governed by the number of monomials of degree $\le d$ in
the variables of the system, and that number grows with $N$. The same six
rows are all \texttt{inconsistent=true}: their ideal is the unit ideal. So
the \texttt{zero\_dim} flag on a random target refutes a decomposition and
says nothing about the geometry of a successful one. The planted toy rows at
$N=8,12$ are consistent, with finite bases of 28 and 49 elements, and are the
positive control. The README's reading is that the useful theoretical target
for PKM is now a lower bound on the last fall or closure degree, or on the
matrix width, and not another first-fall plot; the $N=20$ cell is projected
at 240--380\,GB and is not to be launched to reconfirm the trend. The
solving-degree data of the Kummer towers themselves are in
\S\ref{sec:fd-pkm} and \S\ref{sec:pa-pkm}.

\paragraph{The follow-up.} The same day's follow-up
(\art{research/ic_weak_assumptions_20261102/FOLLOWUP.md}) added an
arbitrary-precision mirror of the tensor route's D.4 admission guards and a
finite-field unit-explanation proof kernel for the arity-three canonical
representation, which it describes as validated implementation artefacts.
It took no scientific run. Its status line per lane: the tensor route's
predicate cost and aggregate complexity remain uncertified; the PKM
closure-dual target is an exact nullspace certificate but no growing
certificate has been found; the Wagner bucket statistic $W$ has no frozen
native producer; for xedni no theorem comparing the selected height with the
global minimal discriminant is proved.

\section{Six weak points measured on Weil-descent systems}\label{sec:hd-weakpoints}

The obligations were then measured, as exact counts, on one family where
every one of them can be computed: Gaudry--Diem descent for
$E:y^2=x^3+ax+b$ with random $a,b\in\FF_{p^n}$, $n\in\{2,3\}$, factor base
$\{P:x(P)\in\FF_p\}$, groups up to $2^{28}$ for the complete solves
(\art{research/ic_weakpoints_20261102/README.md}; protocol frozen before
data, \texttt{PROTOCOL.md} with Amendments 1--3a; native Rust tool, standard
library only, one 4-vCPU container). The README states that it is not an
ecbench session, not a leaderboard row and not admitted by the validation
gate, and reports no wall time. Table~\ref{tab:hd-weakpoints} lists the six
points in the README's numbering.

\begin{table}[tbp]
\centering\small
\begin{tabularx}{\linewidth}{@{}>{\raggedright\arraybackslash}p{0.17\linewidth}>{\raggedright\arraybackslash}X>{\raggedright\arraybackslash}p{0.2\linewidth}@{}}
\toprule
Weak point & Measured (exact counts) & Producer's reading \\
\midrule
W1 solving degree, matrix growth & $n=3$: solving degree 10 in 88/88 systems, $p=31$ to $1{,}073{,}741{,}789$ ($2^5$--$2^{30}$); Macaulay matrix $252\times286$ at that degree, $c(3)=1.03$--$1.07\cdot10^6$ $\FF_p$ multiplications (dense, an upper bound for F4); no first fall in 87/88. $n=2$: degree 3, $6\times10$, about 87 multiplications & a $p$-independent constant $c(n)$; its size, not its growth, is the question \\
\addlinespace
W2 non-generic filter & not addressable on this family (the subfield base is the structure; $\FF_p$ is untouched) & open; the prime-field question \\
\addlinespace
W3 yield mean and variance & 20 curves, 2{,}000 or 6{,}000 targets: mean$/\lambda_{\rm pred}=0.88$--$1.19$ ($\lambda=F^3/6N$); $P(0)=e^{-\rm mean}$ to three decimals; dispersion interval contains 1 on 19/20 curves; the two heavy tails are subfield targets $x_R\in\FF_p$ & Poisson; drop subfield targets \\
\addlinespace
W4 relation rank & 15/15 complete solves return the correct $\log Q$; rank shortfall 0 on 7 curves, 1--3 on 7, 6 on one; cofactor scaling, duplicated rows and dropped signs as controls & not a blocker at this scope \\
\addlinespace
W5 geometry & zero-dimensional 176/176; quotient dimension $64$ (Bézout) on 87/88, 4 at $n=2$; at $p\le127$, 24 targets: 32 rational solutions, of which 5 genuine and 27 non-split ($84\%$), 2 singular, no positive-dimensional component & benign; about five rejections per relation \\
\addlinespace
W6 end-to-end cost against matched rho & complete toy solves, $\ell=2^{12}$--$2^{25}$; IC/rho (negation map, analytic $\sqrt{\pi\ell/4}$): $0.29$--$3.4$ at $w=0$, $2.8\cdot10^4$--$3.2\cdot10^5$ at $w=0.1$, $1.5\cdot10^5$--$1.7\cdot10^6$ at $w=1$; rho ahead on 7 curves, 8 unresolved, IC ahead on none & 4--5 orders of magnitude remaining at $w\ge0.1$ \\
\bottomrule
\end{tabularx}
\caption[]{The six weak points on Gaudry--Diem descent over $\FF_{p^n}$,
$n\in\{2,3\}$, as reported in
\art{research/ic_weakpoints_20261102/README.md}. Unverified producer
figures; not an ecbench session; not admitted by the validation gate. The
eight unresolved curves are unresolved only because $w=0$ makes the Gröbner
work free.}
\label{tab:hd-weakpoints}
\end{table}

The README's reading is that on this family W3, W4 and W5 are cleared, with
cheap filters for their failure modes (subfield targets, duplicate rows,
spurious solutions), that W1 is a constant and not a growth, and that the
factor-base size is a seventh, required lever. Its derived figures, which
are models and not measurements:
\begin{itemize}
\item \emph{No large primes.} With $\tau=3.50$ trials per $p$,
  $c(3)=1.06\cdot10^6$ and $3.97$ nonzeros per row, the best factor-base
  fraction gives IC/rho $\to684$ at $w=0.1$ and $3{,}761$ at $w=1$, constant
  for groups of $2^{60}$ and more. Without large primes $n=3$ balances at
  $p^{3/2}$, rho's exponent, so the gap is a constant proportional to
  $\sqrt{w\,c(3)}$.
\item \emph{Double large primes.} At $|{\rm FB}_r|\approx p^r/2$ the fitted
  exponent of trials to a square system is $1.27\pm0.04$ at $r=2/3$
  (theory $4/3$, rho $1.5$), from $p\le2039$; $\log Q$ was correct in
  105/105 reached cells. Extrapolated with $c(3)$ and the measured row
  weights, 2LP reaches IC/rho of 58, 7.7 and about 1 at $2^{96}$, $2^{132}$
  and $2^{180}$ at $w=0.1$, and 319, 42 and 5.3 at $2^{96}$, $2^{132}$ and
  $2^{180}$ at $w=1$, where it does not cross 1 by $2^{240}$. The band
  therefore ranks IC ahead nowhere up to $2^{240}$: above $2^{180}$ the order
  flips between $w=0.1$ and $w=1$. The README calls this crossover
  optimistic: the fits are extended by up to 70 bits, and the 2LP slopes at
  $r=3/4$ and $5/6$ are probably pre-asymptotic.
\end{itemize}
For the theory of this chapter the result is that, on the one family where
the theorems of Gaudry and Diem apply, the obligations other than W1 and W6
hold at toy size, and the distance that remains is a constant set by the
cost of one decomposition. For prime fields, where W2 is the open
obligation, the family says nothing.

\section{Petit--Quisquater Assumption~2}\label{sec:hd-pq2}

Petit and Quisquater~\cite{petit2012} derived their heuristic
$2^{c\,n^{2/3}\log n}$ for binary curves (\S\ref{sec:theory}) from an
assumption that the corpus records as
\texttt{HEUR-GB-PQ-ASSUMPTION-2}: the degree of regularity, read as the
maximal degree a Gröbner solver reaches, of a Weil-descended summation
system equals its first fall degree up to $o(1)$. At $m=2$ the first fall
degree of the descended $\Sem3$ system is 2 for every $n$ (Kosters and
Yeo~\cite{kosters2015notes}, Corollary~4.11; Corollary~\ref{cor:ffd2} in
our Caminata--Gorla convention), so the literal reading predicts solving
degree 2. Kousidis and Wiemers~\cite{kousidis2019} read the assumption as
$D_{\rm reg}=m^2-m+1+o(1)$, which is $3+o(1)$ at $m=2$. Degrees are
integers, so ``$+o(1)$'' means eventual equality, and a lower bound that
exceeds the predicted value for infinitely many $n$ refutes that reading
for the family.

\begin{theorem}[Subspace criterion; kernel-checked]\label{thm:hd-subspace}
Let $K$ be a field, $I\subseteq K[x]$ an ideal with $1\notin I$, and $\prec$
a degree-compatible monomial order. Suppose $I$ contains $k$ polynomials
whose leading monomials are $k$ distinct variables. If a Gröbner basis of
$I$ lies in a finite-dimensional subspace $S$, then
$k\le\dim_K\bigl(S\cap K[x]_{\le1}\bigr)$. In particular, if $S$ is the span
of the Macaulay rows of degree $\le d$ and $\dim(S\cap K[x]_{\le1})<k$, no
Gröbner basis of $I$ lies in $S$.
\end{theorem}

The Lean declarations are
\texttt{SolvingDegreeBound.linear\_count\_of\_gb\_in\_subspace} and
\texttt{SolvingDegreeBound.no\_gb\_in\_rows\_of\_low\_lt}
(\art{research/ecdlp_weakness_lean_20261104/lean/SubspaceCriterion.lean}).
Its degree-2 special case, \texttt{solving\_degree\_ge\_three} in
\texttt{SolvingDegreeBound.lean}, adds the rank--nullity count
$\dim(\operatorname{span}F\cap K[x]_{\le1})+\operatorname{rank}(\operatorname{top}_2F)\le\abs F$:
for generators of degree 2, field equations included, if
$k+\operatorname{rank}(\operatorname{top}_2F)>\abs F$ then the solving degree
is at least 3. A non-vacuity instance over $\FF_2$ is
$F=\{x_0^2+x_0,\,x_1^2+x_1,\,x_0x_1+x_0+x_1\}$, with $k=2$, rank 3 and
$\abs F=3$. A further declaration, \texttt{not\_tendsto\_of\_frequently\_gap},
proves that an integer gap of at least 1 infinitely often is not $o(1)$.

Write $d_{\rm lin}$ for the least $d$ at which the degree-$\le d$ Macaulay
row space contains $k$ independent polynomials of degree $\le1$, and
$d_{\rm iter}$ for the same quantity computed on the iterated degree-$d$
closure (every multiple of a newly found polynomial that stays within degree
$d$ is added until the rank stops growing). By Theorem~\ref{thm:hd-subspace},
$d_{\rm lin}$ is a lower bound on the Caminata--Gorla one-shot solving degree
(Proposition~\ref{prop:bridges} relates it to the last fall degree), and
$d_{\rm iter}$ is a lower bound on F4's maximal step degree under the
containment property that F4's output at maximal step degree $D$ lies in the
degree-$D$ closure. That property is argued in the census README and is not
formalised. The number $k$ of linear consequences is computed as
$N-\operatorname{affdim}(Z)$ from the enumerated solution set $Z$; the
Boolean Nullstellensatz step behind it is kernel-checked only for $\abs Z=1$
(\texttt{BooleanNullstellensatz.linear\_mem\_of\_unique\_zero}), and the
general affine-span case is not formalised.

\paragraph{The census.} \texttt{pq2\_rank}
(\art{research/ecdlp_weakness_census_20261104/pq2_rank/README.md}; Rust, no
dependencies, seed \texttt{0x5eed20261104}) computes these quantities
exactly over $\FF_2$, with no Gröbner run, on the curve
$y^2+xy=x^3+ax^2+b$ with a random $\lceil n/2\rceil$-dimensional factor-base
subspace and a target with a non-empty solution set. Each instance checks,
on 20 random assignments, that the descended equations agree with a direct
$\Sem3$ evaluation. The main sweep was stopped during $n=47$, before the
binary wrote its own JSON; a script copied the log lines and recomputes
nothing. $n=48$--$56$ were not reached. Table~\ref{tab:hd-pq2} gives the
results.

\begin{table}[tbp]
\centering\small
\begin{tabular}{@{}lrlll@{}}
\toprule
$n$ & instances & iterated degree 3 reaches $k$ & $d_{\rm iter}$ & $d_{\rm lin}$ \\
\midrule
8--14 & 21 & 21 & 3 & 4 (20), $>4$ (1) \\
15 & 3 & 0 (\texttt{iter3}$\,=1$) & 4 & $>4$ \\
16 & 3 & 3 & 3 & $>4$ \\
17--20 & 12 & 0 (\texttt{iter3}$\,=1$) & $>3$ & $>4$ \\
21--44 & 72 & 0 (\texttt{iter3}$\,=1$) & $>3$ & $>3$ \\
45--47 & 8 & not run & -- & $>3$ \\
\midrule
17--24 (fresh) & 24 & 0 & 4 (24/24) & \\
\bottomrule
\end{tabular}
\caption[]{\texttt{pq2\_rank}: exact degree lower bounds for the $m=2$
descended $\Sem3$ system over $\FF_{2^n}$, 119 instances in the main sweep.
The last row is two separate degree-4 sweeps at $n=17$--20 and 21--24 on
fresh instances. At every $n$, \texttt{low2}$\,=1$; \texttt{low3}$\,=1$ for
every $n\ge9$. Unverified census figures. Source:
\art{research/ecdlp_weakness_census_20261104/pq2_rank/README.md}.}
\label{tab:hd-pq2}
\end{table}

\begin{itemize}
\item \emph{Literal Assumption~2 fails on all 119 instances:}
  $d_{\rm iter}\ge3$ everywhere, against the predicted first fall degree 2.
\item \emph{The Kousidis--Wiemers value is not reached} at any of the 87
  instances with $n=15$ or $17\le n\le44$: the iterated degree-3 closure
  reaches only one of the $k$ linear consequences. On the 24 fresh instances
  at $n=17$--24 the iterated degree is exactly 4.
\item \emph{Cross-checks.} A pure-Python closure by naive saturation, sharing
  no code with the Rust one, agrees on five dumped instances at $n=14$--18,
  at degrees 3 and 4. A same-seed rerun for $n=8$--40 reproduces the log
  line for line; it also gives $\abs{F'}/\abs V$ between $0.84$ and $1.31$
  (mean $1.021$) over 99 instances, and $0.997$--$1.004$ for $n\ge30$.
\end{itemize}

What this does and does not settle. A finite range cannot refute an $o(1)$
statement. It refutes the literal reading on every instance computed, and
shows that the asymptotic value 3 is not reached anywhere from $n=17$ to 44.
The bound on F4's step degree rests on the unformalised containment
property; the bound on the one-shot solving degree is a kernel-checked
criterion applied to an unverified computation. The measured degree sits
above the value the subexponential binary heuristics assume, which removes
support from them; it does not support a subexponential algorithm. The
corpus node that records the census, \texttt{MEAS-PQ2-DLIN-CENSUS},
obstructs \texttt{HEUR-GB-PQ-ASSUMPTION-2}, which remains a heuristic in the
corpus. The certified binary instances of Theorem~\ref{thm:lfd5} reach
$d_{\tilde F}\ge5$ at $n=33$--45 on random subspaces of dimension
$\lceil n/2\rceil$; together the two give a lower bound of 4 at $n=17$--24 by
census and of 5 at $n\ge33$ by certificate. What remains open is the same as
in \S\ref{sec:fd-summary}: a proof that the bound grows with $n$.

\section{Semaev's Assumption~1}\label{sec:hd-semaev}

Semaev~\cite{semaev2015} states two degree assumptions for his chained
systems over $\FF_{2^n}$. Assumption~1 (\S4.5) asserts $d_{F4}\le4$ for any
subspace of dimension $k=\lceil n/m\rceil$ and any $2\le t\le m$, supported,
as the corpus records, by 2{,}300 Magma systems with $t=m$, $n\le21$,
$m\le6$. Corollary~\ref{cor:assump} refutes it as stated, in the cells
$(35,2,2,18)$, $(45,2,2,23)$ with Semaev's own default subspace, and
$(19,3,3,7)$, $(21,3,3,7)$ inside his tables' range. The revision of the
paper (ePrint 2015/310, 10 April 2015) replaces it with a growth condition on
the regularity degree that the corpus records as
$o(\sqrt{n/\ln n})$ (\texttt{HEUR-SP-SEMAEV15-WEAK-REGULARITY}), sufficient
for the $2^{c\sqrt{n\ln n}}$ asymptotic and not for the FIPS-curve
conclusions. Nothing the programme proved bears on that weaker condition:
only a growing lower bound would.

His first-fall-degree assumption $d_{F4}\le d_{\rm ff}$ (\S4.4), applied to
his own $t=2$ cells with his statement that the first fall degree ``is always
4'', was tested on 200 instances of the ten cells of his Table~2
(\art{research/semaev_t2_trace_20261102/REPORT.md}; protocol frozen; exact
standard-library arithmetic). The first fall degree is 2 on 200/200, each
with an explicit trace certificate: with $\lambda=z^{-2}$,
\[
  \operatorname{Tr}\bigl(z^{-2}\Sem3(x_1,x_2,z)\bigr)
  =\operatorname{Tr}(x_1)+\operatorname{Tr}(x_2)+\operatorname{Tr}(Bz^{-2}),
\]
a constant combination of the coordinate equations of degree $\le1$, the
Kosters--Yeo identity in Semaev's coordinates. None of the 200 systems has a
solution in $V\times V$ (brute force). The degree at which a Macaulay
computation proves this is 2 on 154 instances and 3 on 46, never 4. So the
statement ``first fall degree always 4'' is refuted, while the inference that
these rows refute $d_{F4}\le d_{\rm ff}$ splits: the sharp assumption fails on
46/200, where 3 exceeds 2, and Semaev's reported $d_{F4}=4$ exceeds what the
systems require. The corpus inference that had read the rows as a refutation
everywhere is corrected in Chapter~\ref{ch:corrections}. The
\texttt{RT-SP-SEMAEV-A1-N45} target of the triage, an exact F4 run in the
$n=45$ cell, is still listed as open.

\section{Nagao's Lemma~4: four cells and three repairs}\label{sec:hd-nagao}

Nagao's ePrint 2015/984~\cite{nagao2015}, which the corpus records as
``polynomial-time binary ECDLP under FFDA'' (status conditional), bounds the
first fall degree $d_F$ by a ``fake'' first fall degree $d'_F$ computed
modulo the field equations (Lemma~4: $d_F\le d'_F$). Remark~\ref{rem:defs}
gives the equality-reading counterexample over $\FF_2$. The lemma's verdict
depends on two choices: the reading of Definitions~5 and~6 (``equality'':
the maximal degree of the products is exactly $d$; ``literal $\ge$'') and
the placement of the field equations (outside or inside the generating set).
Table~\ref{tab:hd-nagao} lists the four cells and the three candidate repairs
tested after the refutation.

\begin{table}[tbp]
\centering\small
\begin{tabularx}{\linewidth}{@{}>{\raggedright\arraybackslash}p{0.24\linewidth}>{\raggedright\arraybackslash}X>{\raggedright\arraybackslash}p{0.24\linewidth}@{}}
\toprule
Statement & Result & Formal status \\
\midrule
Lemma~4, equality, outside & refuted: $\{X^2Y,XY+X\}$, $d'_F=2<d_F=3$ (Remark~\ref{rem:defs}) & two implementations \\
Lemma~4, equality, inside & refuted & corpus record \\
Lemma~4, literal $\ge$, outside & refuted ($d_F=3$, $d'_F=2$) & corpus record \\
Lemma~4, literal $\ge$, inside & \emph{theorem}: every fake fall is a true fall at the same degree, so $d_F\le d'_F$, over any field & kernel-checked (Mathlib) \\
\addlinespace
A1: $d_F\le f(d'_F)$, outside & refuted: $W_k=\{X^kY,XY+X\}$ has $d'_F=2$, $d_F=k+1$ & $k=3$ kernel-checked; all $k$ by hand proof \\
A2: $d_F\le d'_F+1$, inside & refuted: $\{YZ+1,XZ+X+1\}$, $d'_F=1$, $d_F=3$; 10{,}752 of 521{,}731 pairs have gap 2 & native, two implementations \\
Inside gap, $n\ge2$ & $d_F-d'_F\le3$, tight ($\{XYW+XY+W,XZW+XW+1\}$); unbounded for $n=1$ & key identity kernel-checked in two variables \\
A3: $d'_F\le d_F$ when the reduced generators are independent & survives: 0 counterexamples in 25{,}536 and 322{,}896 independent triples & census only \\
\bottomrule
\end{tabularx}
\caption[]{Nagao's Lemma~4 over $\FF_2$ by reading and placement, and the
candidate repairs. Sources:
\art{docs/ecdlp-hardness-corpus/WEAKNESS_TRIAGE.md} (W4),
\art{research/ecdlp_weakness_lean_20261104/lean/NagaoLiteralInside.lean},
\art{research/lean_attacks_20261102/REPORT.md},
\art{research/inside_gap_n4_20261102/REPORT.md}.}
\label{tab:hd-nagao}
\end{table}

\paragraph{The fourth cell.} \texttt{NagaoLiteralInside.lemma4\_literal\_inside}
proves, for any field, any family $F$ and any finite family $E$ of
``field equations'' placed among the generators, that $d_F\le d'_F$
whenever a fake fall exists. The proof lifts a fake fall: replace the sum by a
least-degree representative modulo $\langle E\rangle$ and pay the difference
with multiples of $E$, which is allowed because $E$ is inside. The file's
comment records why the argument fails outside, or under the equality
reading, where the added multiples can push the maximal degree above $d$. The
reading axis is now complete, refuted in three cells and proved in the
fourth. Which cell Nagao meant is a question of the source text: the triage
lists a fresh PDF-level transcription of Definitions~5 and~6 as owed.

\paragraph{A1: outside, the fake degree bounds nothing.} Modulo the field
equations $X^kY\mapsto XY$, so every $W_k$ has the reduced system of the
original witness and $d'_F=2$. With the field equations outside,
$X^{k-1}(XY+X)+X^kY=X^k$ is a fall at $k+1$, and no fall exists below it, so
$d_F=k+1$: the native table gives $3,4,\ldots,8$ for $k=2,\ldots,7$, and 3
for every $k$ with the field equations inside
(\art{research/lean_attacks_20261102/results/family.txt}).
\texttt{IcNegative.FakeDegreeGap.gap\_two\_witnessW3} kernel-checks $k=3$
($d'_F=2$, $d_F=4$) in core Lean with \texttt{decide +kernel}; the
repository proof gate passed it in run 41399938870 (target
\art{papers/ic-negative-results/lean} PASS, 40 theorems, 0 trusted). The
all-$k$ statement is the hand proof plus the table.

\paragraph{A2 and the inside gap.} The census of all 521{,}731 pairs of
degree $\le2$ in three variables (\art{research/lean_attacks_20261102/results/p3gap.txt})
finds gap 2 inside on 10{,}752 pairs; the witness $\{YZ+1,XZ+X+1\}$ has a
fake fall at 1 and, inside, the true fall
$Z(XZ+X+1)+X(Z^2+Z)=Z$ at 3. Two implementations agree (a Rust census and a
Python brute force over every admissible multiplier tuple); the witness has
three variables and is outside the two-variable Lean encoding. The open
question this left was answered the same day
(\art{research/inside_gap_n4_20261102/REPORT.md}; protocol frozen): for
$n\ge2$ variables, $d_F(F\cup S_{\rm fe})\le4$ for every system, because
$Y^2(X^2+X)+X^2(Y^2+Y)=XY^2+X^2Y$ is a fall at 4 built from the field
equations alone, so the inside gap is at most 3; a gap-3 witness in four
variables was found in a sample of $2{,}000{,}000$ pairs, and for $n=1$ the
family $\{X^k,X^k+1\}$ has gap $k-1$. \texttt{InsideGap.lean} states the
field-equation fall in two variables in core Lean; the gate passed it in
run 41484047855 (44 theorems, 4 new, 0 trusted). The step ``zero multipliers
on further generators leave the fall intact'' is argued, and the
four-variable witnesses are native only. The consequence is that
inside-placement $d_F$ under the equality reading is degenerate: it never
exceeds 4 in two or more variables, through the trivial falls among field
equations that Ding--Hodges and Petit--Quisquater exclude from their first
fall degree. No estimate can take its degree from that quantity without
first removing them.

\paragraph{A3 and Nagao's Proposition~5.} The reverse repair, $d'_F\le d_F$
outside when the reduced generators are linearly independent, has no
counterexample among the triples tested: 266 counterexamples among the
dependent triples of degree $\le2$ in two variables, none among the 25{,}536
independent ones, and none among the 325{,}500 multilinear triples in three
variables (\art{research/lean_attacks_20261102/REPORT.md}). Nothing here
touches Nagao's Proposition~5 or the paper's own applications, which check
unreduced degrees instance by instance; only a bound through $d'_F$ alone is
ruled out, and only outside.

\section{Relation rank and the left kernel}\label{sec:hd-rank}

The corpus node \texttt{HEUR-RELATION-RANK} conditions sparse linear algebra
and, through it, marked three externally proved results as resting on a
heuristic: Theorem~1 of Gaudry, Thomé, Thériault and
Diem~\cite{gtthd2007}, Diem's small-genus theorem~\cite{diem2006} and the
GHS Weil-descent attack~\cite{ghs2002}. These are 4 of the 17 status
inflations of \S\ref{sec:hd-corpus}. Galbraith and Gaudry~\cite{galbraith2016}
(\S6.2) already note that target-dependent relations are easier to analyse
rigorously than relations reused across targets. The triage wrote the
argument out (W2) and the Lean round of 4~November formalised its finite core.

\begin{theorem}[Left-kernel relations; kernel-checked]\label{thm:hd-leftkernel}
Let $K$ be a field and $M\in K^{n\times B}$ a relation matrix with
$a_i+b_ix=\sum_jM_{ij}f_j$ for $i=1,\ldots,n$, where $f_j$ are the
factor-base logarithms and $x$ the unknown logarithm.
\begin{enumerate}
\item If $n>B$, a nonzero $\lambda\in K^n$ with $\lambda M=0$ exists, for
  every $M$.
\item For any such $\lambda$,
  $\sum_i\lambda_ia_i+x\sum_i\lambda_ib_i=0$; hence
  $x=-\sum\lambda_ia_i/\sum\lambda_ib_i$ whenever $\sum\lambda_ib_i\ne0$.
\item Let $K$ be finite with $q$ elements and let $\Lambda$ be any
  function of the vector of relation logarithms $\ell=(a_i+b_ix)_i$ alone
  whose value is always a nonzero vector. Exactly $q^{2n-1}$ of the $q^{2n}$ pairs
  $(a,b)$ give $\Lambda(\ell)\cdot b=0$: the failure density is exactly $1/q$.
\end{enumerate}
\end{theorem}

The declarations are \texttt{LeftKernel.exists\_left\_kernel},
\texttt{left\_kernel\_identity}, \texttt{left\_kernel\_recovers} and
\texttt{card\_bad\_pairs} (stated without subtraction as
$\abs{\rm bad}\cdot q=q^{2n}$) in
\art{research/ecdlp_weakness_lean_20261104/lean/LeftKernel.lean}. Two steps
are not formalised: that a decomposition of $R_i=a_iP+b_iQ$ gives the row
$M_i$ (true by the definition of the factor-base logarithms), and the
reduction from uniform $(a_i,b_i)$ to the counting statement, which is the
uniform-measure statement itself. The consequence for the corpus: for a
single target, with relations drawn from target-dependent points, index
calculus needs no rank condition at all, and the four inflation rows are edge
errors in the graph, not fragile theorems. Removing the heuristic
\emph{strengthens} the closures that rest on it, such as the GHS and
two-summand closures. The triage's own toy check (12 seeds at $p=10007$,
every matrix 2--7 short of full rank, left-kernel recovery correct 12/12) was
run in scratch space and is labelled there as an illustration, not evidence.

The heuristic stays load-bearing where relations are reused across targets:
right-kernel amortisation, multi-target batches, and PKM's surplus of
$\Delta=10$ relations (\texttt{HEUR-PF-PKM-RANK-DELTA}).
\texttt{rank\_deficit}
(\art{research/ecdlp_weakness_census_20261104/rank_deficit/README.md};
results in \texttt{results/tables.md}) measures that regime exactly: ten
prime-order curves, two per $p\in\{1009,2003,4001,8009,16001\}$, small-$x$
bases with $B$ closest to $\sqrt r$ ($B=32$ to 126), rows $aG=P_1+P_2$, five
samples each.
\begin{itemize}
\item With negation pairs merged into $B/2$ unknowns, $\Delta=10$ gave full
  rank in all five samples on 2 of the 10 curves; on the other eight the
  largest deficit at $\Delta=10$ is 1--3. At $\Delta=40$ four curves still
  have a deficit in some sample, at most 2. On four curves no tested
  $\Delta\le40$ gives full rank in all five samples.
\item On the unmerged $B$ columns no tested $\Delta\le40$ gives deficit 0 in
  all five samples on any curve.
\item Free relations $P_1+P_2=P_3$ among factor-base points span only 3 to
  15 of the $B/2$ merged unknowns.
\end{itemize}
At these sizes PKM's ``$\Delta=10$ suffices'' is contradicted
(\texttt{MEAS-RANK-DEFICIT-CENSUS} obstructs \texttt{HEUR-PF-PKM-RANK-DELTA}).
The census README also states that the unmerged deficit is ``at least about
$B/4$ at every $\Delta$''; its own table does not show this (at $\Delta=40$
the minimum unmerged deficit over five samples is 0 or 1 on seven curves), so
that sentence is not used here. The measured complete solves agree with the
theorem: on the descent family 15/15 solves returned the correct logarithm
with rank shortfalls up to 6 (\S\ref{sec:hd-weakpoints}), and the 842 prime
$m=2$ rows of the reanalysis report rank 842, a correctness control and not
a rank law, since each episode stopped before a full matrix was collected
(\art{research/ic_weak_assumptions_20261102/evidence-summary.json}).

\section{Yield: the mean is a theorem, the shape is a heuristic}\label{sec:hd-yield}

The corpus node \texttt{HEUR-DECOMPOSITION-PROBABILITY} has the largest
impact on the frontier of the attack graph (13). Its mean, $B^m/r$, follows
by counting: \texttt{MeanCount.sum\_card\_reps} kernel-checks that the
ordered representation counts over all targets sum to $\abs F^m$, the
ordered form of Proposition~\ref{prop:conservation}. Only the shape is
heuristic: Poisson zero mass, variance, independence between targets,
rationality and visibility to the solver. The triage (W3) splits the node in
two. Over $\FF_{q^n}$ with $n$ fixed, Diem's rigorous results bypass the
heuristic. Over prime fields it is genuinely heuristic. Under the
left-kernel formulation the cost needs only the fraction of points that
decompose, a set size, exact on small fields.

\paragraph{Prime field, $m=2$.} \texttt{pf\_yield}
(\art{research/ecdlp_weakness_census_20261104/pf_yield/results/table.md})
enumerates every target on 18 prime-order curves, $p$ from 2003 to 80021,
with three base shapes (small $x$; random; random closed under negation) at
$\mu=B^2/r$ near $0.5$, 1 and 2: 162 rows. With ordered pairs the
decomposable fraction is $0.55$--$0.78$ of $1-e^{-\mu}$. With unordered
pairs $U(R)$, of mean $\mu/2$, $P(U\ge1)$ matches $1-e^{-\overline U}$: the
mean absolute deviation of the ratio is $0.011$, $0.007$ and $0.011$ for the
three shapes. The variance-to-mean ratio of $U$ is not Poisson for the bases
index calculus uses: averaged over 54 rows it is $1.455$ for small $x$ and
$1.447$ for the random negation-closed base, against $0.981$ (range
$0.91$--$1.01$) for the random base that is not closed under negation. The
overdispersion follows closure under negation, through $U(R)=U(-R)$, which
couples targets in pairs, and not the small-$x$ shape. The heuristic's zero
class is right once the pairing convention is stated; its variance is off by
a factor of about $1.45$.

\paragraph{The descent family and the $1/n!$ proportion.} On the
Gaudry--Diem family of \S\ref{sec:hd-weakpoints} the yield is Poisson:
mean$/\lambda_{\rm pred}=0.88$--$1.19$ and $P(0)=e^{-\rm mean}$, the heavy
tails being subfield targets. \texttt{ext\_decomp}
(\art{research/ecdlp_weakness_census_20261104/ext_decomp/results/table_primary.md})
measures the $n$-fold decomposable fraction over $\FF_{q^n}$ with factor base
$x\in\FF_q$ against the survey value $1/n!$: $0.350$--$0.413$ at $n=2$
($q=101,211,401$) and $0.101$--$0.165$ at $n=3$ ($q=23,41,61$), that is
$0.70$--$0.83$ and $0.61$--$0.99$ of $1/n!$ (the README's summary,
$0.70$--$0.99$ for both, overstates the lower end at $n=3$). The $n=4$ rows
have factor bases of 10--18 points and scatter widely. For control curves
defined over $\FF_q$ the premise fails: the factor base is $E(\FF_q)$ or lies
in a small subgroup, and the fraction is $0.0003$--$0.002$ at $n=3$. The
$1/n!$ is the mean of Proposition~\ref{prop:conservation} and carries the
hypothesis that the curve is not defined over the subfield.

\paragraph{A withdrawn closure.} The triage first argued by a dimension
count that the trace pull-back (base-change $E/\FF_p$ to $\FF_{p^n}$,
decompose over points with $x$ in an $\FF_p$-subspace $V$, pull back by the
trace) cannot give a small factor base, and withdrew the argument on
4~November: the count bounds $\abs{F'}$, not $\abs{\operatorname{Tr}(F')}$.
\texttt{trace\_pullback}
(\art{research/ecdlp_weakness_census_20261104/trace_pullback/results/summary_table.md})
computes the trace image of every one-dimensional subspace for ten curves,
$p=103$ to 2003, over $\FF_{p^2}$: the smallest image is $0.319$--$0.389$ of
$\#E(\FF_p)$, attained at $V=\FF_p\cdot i$ on every curve, and two-point
decomposition over it covers every target. The lead
\texttt{GAP-PF-TRACE-PULLBACK-FACTOR-BASE} is narrowed, not closed: higher
extension degrees and higher-dimensional $V$ are not covered.

\section{Waiting times: marginals suffice, hazards are not implied}\label{sec:hd-hazard}

The Jacobson--Koblitz--Silverman--Stein--Teske obstruction to xedni,
audited in Lean on 1~November (\art{research/jksst_xedni_audit_20261101/};
see Chapter~\ref{ch:corrections} for the correction of the treatise's earlier
``closed'' wording), derives a linear expected waiting time from a
conditional-hazard recurrence $(1-\varepsilon)s(n)\le s(n+1)$ on the survival
masses $s(n)=P(\tau>n)$, and listed the recurrence for repeated adaptive
trials as open. The triage (W7) promoted that free-text assumption to a
node, \texttt{ASM-XEDNI-HAZARD-RECURRENCE}.

\begin{theorem}[Marginals suffice; kernel-checked]\label{thm:hd-hazard}
Let $\mu$ be a probability mass on a finite set, $\tau$ a stopping time, and
$S_t$ the event that trial $t$ succeeds, with $\tau\le n$ implying some
$S_t$, $1\le t\le n$. Trials may be correlated or adaptive.
\begin{enumerate}
\item If $\mu(S_t)\le\varepsilon$ for every $t\le N$, then
  $s(n)\ge1-n\varepsilon$ for all $n\le N$.
\item If moreover $N\varepsilon\le1/2$, then
  $\mathbb E[\min(\tau,N)]\ge N/2$.
\item The converse inference fails: on three equally likely outcomes, outcome
  $k$ succeeding at trial $k+1$, every trial has marginal success mass at
  most $1/3$, the recurrence with $\varepsilon=1/3$ fails at $n=1$
  ($\tfrac23\cdot\tfrac23=\tfrac49>\tfrac13=s(2)$), and still
  $\mathbb E[\min(\tau,3)]=2$.
\end{enumerate}
\end{theorem}

The declarations are \texttt{survival\_ge\_of\_marginals},
\path{truncated_expectation_ge_half_horizon_of_marginals},
\texttt{marginals\_U}, \texttt{marginals\_do\_not\_give\_hazard} and
\texttt{wait\_U} in
\art{research/ecdlp_weakness_lean_20261104/lean/HazardSeparation.lean}.
Part~2 is the conclusion of the JKSST audit's
\texttt{truncated\_expectation\_ge\_half\_horizon} without its hazard
hypothesis; part~3 refutes the corpus inference
\texttt{CLAIM-MARGINAL-IMPLIES-HAZARD}. The open bridge for xedni is
therefore a per-trial marginal bound for the adaptive trial actually run
(\texttt{ASM-XEDNI-TRIAL-MARGINAL}: each trial succeeds with marginal
probability at most the one-shot bound $C_0/p$), a weaker and more natural
premise than a history-uniform hazard bound. This strengthens the
obstruction; it does not reopen xedni, and it leaves the other conditional
links of the JKSST chain (heights and discriminants, Lang's conjecture,
selected lifts, useful relations, cofactors) as they were.

\section{Generic bounds and what they cover}\label{sec:hd-generic}

The non-generic-interface obligation exists because a generic algorithm
cannot beat $\sqrt r$ (\S\ref{sec:gf-shoup}). Two Lean results this autumn
fix the scope of that statement.

\paragraph{Shoup's bound is relative to random encodings.}
\path{TransparentEncoding.dlog_identity_encoding} and
\path{zero_operation_dlog}
(\art{research/ecdlp_weakness_lean_20261104/lean/TransparentEncoding.lean})
prove that with the identity encoding of $\ZZ/r$ the logarithm is the
encoding, read with zero group operations. This refutes the corpus claim
\texttt{CLAIM-SHOUP-ALL-ALGORITHMS} (``every algorithm needs
$\Omega(\sqrt r)$'') and qualifies the bound and ``BSGS is optimal'' as
generic-model statements. It is a scope statement, not an attack: it is the
reason every subexponential claim must name an encoding property, and the
reason the obligation of Table~\ref{tab:hd-chain} asks for a coordinate
predicate that no generic table reproduces. The generic-interface audit of
1~November (\art{docs/ecdlp-hardness-corpus/GENERIC_INTERFACE_AUDIT.md};
an internal argument, not a reviewed theorem) sorts the programme's oracles:
subtract, meet-in-the-middle, $k$-sum and pair-enumeration oracles are
generic table algorithms once their set-up is charged; the root-finding
oracles return the same relations as group subtraction, with
representation-specific set-up; Frobenius-orbit methods need an augmented
interface with a cheap known-scalar automorphism; and F4/F5/SAT, Weil-descent,
PKM-tower and small-root lattice oracles are coordinate-algebraic, so Shoup's
and the Corrigan-Gibbs--Kogan~\cite{cgk2018} bounds do not exclude them. Only
the last group can discharge the obligation. The weak-assumption reanalysis
credits the PKM/F4 and compact Boolean predicates with a genuinely
coordinate-algebraic interface; none of them yet has a solver cost that
discharges the obligations after it (\S\ref{sec:hd-chain},
\S\ref{sec:pg-gate}).

\paragraph{Rho is conductor-blind in distribution.} The conductor-gap
programme asked whether the depth of $\operatorname{End}(E)$ in an isogeny
volcano changes the cost of rho. The theorem that answers it in distribution
is \art{research/conductor_gap_rho_large_sample_20261018/lean/RepresentationHashInvariance.lean}.

\begin{theorem}[Representation-hash invariance; checked by the producer]\label{thm:hd-rhoinv}
Let $\varphi:G\to H$ be a bijection compatible with binary operations on
$G$ and $H$. For an $r$-adding walk $w=(h,c)$ on $G$, with partition
$h:G\to\{0,\ldots,r-1\}$, multipliers $c$ and step $x\mapsto x\cdot c(h(x))$,
let $\varphi_*w=(h\circ\varphi^{-1},\varphi\circ c)$. Then the trajectory of
$\varphi_*w$ from $\varphi(x)$ is $\varphi$ of the trajectory of $w$ from
$x$; the rho length (first repeat, capped at the budget) is the same at every
budget; and $w\mapsto\varphi_*w$ is a bijection of walk spaces.
\end{theorem}

No group axiom is used. Read outside Lean, the push-forward of the uniform
distribution on walks under a bijection is uniform, so over a uniformly
random partition the cost distribution of rho is the same on the crater and
on the floor, and on every curve of an isogeny class of prime order (abstract
groups of equal prime order are isomorphic). The conductor, the
$j$-invariant or any invariant that is not an invariant of the group can
enter only through the one fixed hash an implementation uses; whether
ecbench's fixed partition does so is an empirical question, and the
measurements are in \S\ref{sec:rv-isogeny}. Formal status: 11 theorems in
core Lean~4.34.1 without Mathlib or a Lake project; the committed
\texttt{build.log} of 18~October reports \texttt{[propext, Quot.sound]} for
the three audited declarations and no \texttt{sorryAx}. The file is not
among the targets of proof-gate run 41399938870. Nothing about curves,
isogenies, endomorphism rings or conductors is formalised, and the file says
nothing about index calculus or about the cost of computing the connecting
isogeny. It complements \texttt{GenericInvariance.lean} of
\art{research/conductor_gap_cm_pairs_20261010/lean}, which covers only
algorithms whose control flow depends on equality answers (the gate run above
passes that project with 60 theorems); see also \S\ref{sec:mr-lean}.

\section{The hardness corpus}\label{sec:hd-corpus}

\art{docs/ecdlp-hardness-corpus/} is a typed claim graph of what supports
elliptic-curve discrete-log security: definitions, assumptions, reductions,
algorithms, model lower bounds, conditional attacks, heuristics,
measurements, obstructions and open obligations, each with its model and
scope, and every edge typed (\texttt{supports}, \texttt{conditions},
\texttt{obstructs}, \texttt{qualifies}, \ldots). The triage of
3~November (\art{docs/ecdlp-hardness-corpus/WEAKNESS_TRIAGE.md}) read the
assembled version of 1~November, 583 claims, 950 edges and 236 sources, and
marks its own findings \texttt{[pcc:unverified]}. The generated
\texttt{stats.tsv} on the current main branch counts 608 nodes, 32 of them
refuted and 68 \texttt{proved\_internal\_unreviewed}; the corpus README's
``current snapshot'' still gives an older 525. Counts below are the
triage's.

\begin{table}[tbp]
\centering\small
\begin{tabularx}{\linewidth}{@{}>{\raggedright\arraybackslash}p{0.27\linewidth}X r@{}}
\toprule
Check & What it finds & Count \\
\midrule
Status inflation & proved, algorithm or measured claims whose support closure reaches an open, heuristic or assumed claim & 17 \\
Unwired dependence & a claim's assumption text invokes a heuristic node with no edge to it & 11 \\
Live obstruction & an established claim obstructs, qualifies or does not cover a claim that still supports others & 49 \\
Route closure on weak support & an \texttt{obstructs} edge whose obstruction rests on an assumed, heuristic or open claim & 16 \\
Hidden shared assumption & the same assumption text on 3 or more nodes, never a node itself & 18 clusters \\
Source contagion & a primary source with a refuted claim and other live claims & 21 \\
Model-scope leap & a generic-model or measured claim supporting a broader one without a model bridge & 9 \\
\bottomrule
\end{tabularx}
\caption[]{Structural checks of the hardness corpus (version of
1~November: 583 claims, 950 edges, 236 sources), as counted by
\texttt{weakness.mjs} and reported in
\art{docs/ecdlp-hardness-corpus/WEAKNESS_TRIAGE.md}. It then ranks 246
attackable claims by impact times breakability and keeps a separate hygiene
queue of 100; the triage states that the scores are a heuristic, not
evidence.}
\label{tab:hd-corpus}
\end{table}

The triage's findings, W1--W7, are the targets of the censuses and Lean
modules of this chapter:
\begin{itemize}
\item \emph{W1: the first-fall pillar is obstructed and still load-bearing.}
  \texttt{HEUR-FIRST-FALL-DEGREE} has four incoming \texttt{obstructs} edges
  from established claims and still conditions seven claims, among them the
  ECC2K-130 $2^{86}$ projection, Nagao's complexity, the PKM algebraic route
  and binary Weil descent. Its siblings \texttt{HEUR-GB-SEMIREG-ASSUMPTION}
  (six \texttt{obstructs} edges) and \texttt{HEUR-GB-PQ-ASSUMPTION-2} (four)
  head the attack queue with it. Because the graph keeps them at
  ``heuristic'', the seven conditional estimates read as one discharged
  condition away from usable, which overstates them
  (\S\ref{sec:hd-pq2}).
\item \emph{W2: relation rank} (\S\ref{sec:hd-rank}).
\item \emph{W3: yield}, two regimes in one node (\S\ref{sec:hd-yield}).
\item \emph{W4: the Lemma~4 refutation rests on a reading choice}, and two
  internal papers disagreed about it; the wording of
  \art{papers/ic-negative-results/} was corrected on 4~November to the
  four-cell table (\S\ref{sec:hd-nagao}).
\item \emph{W5: eleven unwired dependences} understate the impact of the top
  frontier nodes; refuting the heuristic would not reach them.
\item \emph{W7: route closures on weak support.} Xedni is closed twice, both
  softly; lifted index calculus is closed by three heuristics; the naive
  prime-field Semaev method is closed by a lower bound conditional on a
  statistical conjecture validated only to about 25-bit primes; the
  quasi-subfield routes are closed by results resting on the yield
  heuristic. Rank-only closures are not reopenings, by W2.
\end{itemize}
The triage states that W1--W7 change no claim's status by themselves; the
Lean and census rounds added nodes and edges. Its recommended order puts
corpus and proof work first (re-wire the Weil-descent chain to the
left-kernel theorem, add the eleven edges, transcribe Nagao's definitions)
and lists the remaining measurements (YYTK's conjecture past 25 bits,
Semaev's $n=45$ cell, the per-trial xedni marginal) as runs that go through
the experiment route and the validation gate before any number enters the
corpus.

\section{Formal status}\label{sec:hd-lean}

Table~\ref{tab:hd-lean} gathers the formal artefacts of this chapter with
what a kernel has checked and what is only stated or argued. The degree-four
lower bound of the guarded chain (\texttt{GuardedModel}, Theorem~\ref{thm:guarded})
and the other fall-degree identities are in \S\ref{sec:fd-guarded}; the
seven theorems of the Semaev-alternatives audit, which correct an oracle gate
to $s<m-3$, are in \S\ref{sec:pg-gate}.

\begin{table}[tbp]
\centering\small
\begin{tabularx}{\linewidth}{@{}>{\raggedright\arraybackslash}p{0.27\linewidth}>{\raggedright\arraybackslash}X>{\raggedright\arraybackslash}p{0.27\linewidth}@{}}
\toprule
Artefact & Kernel-checked & Not formalised \\
\midrule
\path{ecdlp_weakness_lean_20261104} (Lean and Mathlib v4.34.1): \texttt{LeftKernel}, \texttt{HazardSeparation}, \texttt{SolvingDegreeBound}, \texttt{SubspaceCriterion}, \texttt{MeanCount}, \texttt{BooleanNullstellensatz}, \texttt{NagaoLiteralInside}, \texttt{TransparentEncoding}, \texttt{ProtocolCountermodels} & Theorems~\ref{thm:hd-subspace}, \ref{thm:hd-leftkernel}, \ref{thm:hd-hazard}; the fourth Lemma~4 cell; the mean count; the $\abs Z=1$ Nullstellensatz; 45 \texttt{\#print axioms} lines in \texttt{AxiomAudit.lean} & F4 containment; $k=N-\operatorname{affdim}Z$ for $\abs Z>1$; that PQ's ``maximal degree'' is the Caminata--Gorla solving degree; relation-to-row step \\
\addlinespace
\texttt{IcNegative.FakeDegreeGap} & $W_3$: $d'_F=2$, $d_F=4$ (gate run 41399938870) & $W_k$ for all $k$ (hand proof) \\
\addlinespace
\texttt{IcNegative.InsideGap} & the field-equation fall at 4 in two variables (gate run 41484047855) & zero multipliers on further generators; four-variable witnesses \\
\addlinespace
\path{RepresentationHashInvariance} (core Lean 4.34.1) & Theorem~\ref{thm:hd-rhoinv}, 11 theorems; producer's \texttt{build.log} & curves, isogenies, conductors; the fixed hash \\
\bottomrule
\end{tabularx}
\caption[]{Formal status of the results of this chapter.}
\label{tab:hd-lean}
\end{table}

For the 4~November project the exact status is as follows. Its README says
each module was compiled with \texttt{lake env lean -o} against the Mathlib
v4.34.1 cache, that a full local \texttt{lake build} was not possible for lack
of disk, and that CI's \texttt{lean-proof-gate} runs the full build and the
transitive axiom audit. The triage records that every theorem audited after
round~3 (35) depends only on \texttt{propext}, \texttt{Classical.choice} and
\texttt{Quot.sound}. A search of the committed sources finds no
\texttt{sorry}, \texttt{admit}, \texttt{native\_decide} or \texttt{axiom}.
No proof-gate run identifier is recorded for this project in its directory or
in the corpus fragments, which say the gate ``runs on the landing PR'' or
``on main after merge''; the pcc claim records are described as pending CI
signing. The corpus marks the nodes \texttt{kernel\_checked} with status
\texttt{proved\_internal\_unreviewed}. The theorems are general statements
of linear algebra, counting and finite probability; their bearing on the
literature is through the bridges named in the right column, which are
argued, not proved.

\section{What the literature's subexponential claims now rest on}\label{sec:hd-ledger}

Table~\ref{tab:hd-ledger} sets the claims of \S\ref{sec:theory} and the two
recent prime-field claims against the obligations of
Table~\ref{tab:hd-chain}, with what is proved, refuted, conditional and open.
Statuses in the second column are the corpus's.

\begin{table}[tbp]
\centering\small
\begin{tabularx}{\linewidth}{@{}>{\raggedright\arraybackslash}p{0.2\linewidth}>{\raggedright\arraybackslash}p{0.13\linewidth}X@{}}
\toprule
Claim & Corpus status & What this programme established \\
\midrule
Diem 2011: subexponential ECDLP over $\FF_{q^n}$ for $n$ growing with $\log q$~\cite{diem2011} & proved (external) & Untouched. Its scope, as the corpus records it: for fixed $n$ the bound is exponential, and for $q=2$, $n$ prime it is $e^{O(n^2)}$, far above rho. For the rigorous Weil-descent chain (GTTD, Diem 2006, GHS) the rank step needs no heuristic for single targets (Theorem~\ref{thm:hd-leftkernel}). \\
\addlinespace
Petit--Quisquater 2012: $2^{c\,n^{2/3}\log n}$ for binary curves~\cite{petit2012} & heuristic premise & Literal Assumption~2 fails on 119/119 instances; the $m=2$ Kousidis--Wiemers value 3 is not reached at $n=17$--44; $d_{\tilde F}\ge5$ certified at $n=33$--45. The $o(1)$ form is not refuted: no growing lower bound. \\
\addlinespace
Semaev 2015: $2^{c\sqrt{n\ln n}}$, $c\approx1.69$~\cite{semaev2015} & conditional & Assumption~1 ($d_{F4}\le4$) refuted as stated in four cells; ``first fall degree always 4'' refuted at $t=2$ (200/200 have 2). The revised growth condition is open. \\
\addlinespace
Nagao 2015/984: polynomial time under FFDA~\cite{nagao2015} & conditional & Lemma~4 refuted in three reading--placement cells, proved in the fourth; A1 and the $+1$ form of A2 refuted; inside gap bounded by 3 because inside $d_F\le4$ is degenerate; A3 survives on the triples tested. The complexity claim rests on the first-fall heuristic of W1. \\
\addlinespace
PKM 2016: prime-field algebraic route~\cite{pkm2016} & conditional & $m=4$ solving degrees $6,7,7$ above the threshold 4 and work $\times558{,}304$ from $N=8$ to 16; $\Delta=10$ contradicted at $p\le16001$ in the right-kernel regime; no cheap root-finding algorithm. \\
\addlinespace
Hashed-base Boolean tensor contraction (Akshima et al.\ v0.2) & conditional & The only inspected route with an explicit end-to-end subexponential bound; conditional on the pinned contraction theorem; predicate cost and aggregate complexity uncertified (\S\ref{sec:pt-tensor}). \\
\addlinespace
Xedni and lifted index calculus & obstruction conditional & JKSST is a chain of conditional links; the waiting-time link needs only per-trial marginals (Theorem~\ref{thm:hd-hazard}); the marginal bound for the adaptive trial is open. \\
\bottomrule
\end{tabularx}
\caption[]{The subexponential claims and the obligations they discharge.
Corpus statuses from \art{docs/ecdlp-hardness-corpus/}; the right column
summarises \S\S\ref{sec:hd-chain}--\ref{sec:hd-generic}. Census figures
are unverified; Lean results as in Table~\ref{tab:hd-lean}.}
\label{tab:hd-ledger}
\end{table}

Read by kind:
\begin{itemize}
\item \emph{Proved here.} Five general statements, all kernel-checked and
  unreviewed: the subspace criterion, the left-kernel theorem with its exact
  $1/q$ failure density, the marginal waiting-time bound with its separation,
  the fourth Lemma~4 cell, and representation-hash invariance of rho. None
  is a lower bound on an index-calculus algorithm. Two of them
  \emph{remove} a heuristic: the rank heuristic from the rigorous
  Weil-descent chain, and the hazard recurrence from the xedni obstruction.
\item \emph{Refuted.} Literal readings, all at the level of a stated
  assumption: Semaev's Assumption~1 and his $t=2$ first-fall statement;
  Nagao's Lemma~4 in three cells, and the repairs A1 and A2 ($+1$);
  Petit--Quisquater Assumption~2 read literally, on every instance computed;
  PKM's $\Delta=10$ at toy sizes; the corpus inferences
  \texttt{CLAIM-MARGINAL-IMPLIES-HAZARD} and
  \texttt{CLAIM-SHOUP-ALL-ALGORITHMS}.
\item \emph{Conditional.} Every subexponential claim for binary or prime
  fields remains conditional: on an asymptotic degree assumption
  (Petit--Quisquater's $o(1)$, Semaev's revised growth condition), on a
  first-fall heuristic (Nagao, PKM), or on an external theorem (the tensor
  route).
\item \emph{Open.} A last-fall or solving-degree lower bound that grows with
  $n$ for any presentation (\S\ref{sec:fd-summary}); a per-trial marginal
  bound for adaptive xedni lifts; the trace pull-back beyond
  one-dimensional $V$ over $\FF_{p^2}$; a non-generic prime-field predicate
  with a simulator-or-separation proof; and whether ecbench's fixed rho hash
  sees the conductor.
\end{itemize}
The direction is consistent. Each refutation removes support from a
subexponential claim, and none supplies one; each theorem either confines a
heuristic to a narrower regime or shows a closure to be stronger than its
stated support. On the only family where every obligation was measured, the
Gaudry--Diem descent, the distance to matched rho at toy size is 4--5 orders
of magnitude at $w\ge0.1$, and IC is ahead nowhere in the band.
